% Capitulo: Na prática: projetos de engenharia de dados
\chapter{Na prática: projetos de engenharia de dados}
\label{cap:proj-engenharia}

Projetos de engenharia de dados têm o problema oposto ao dos projetos de visualização: o
resultado é pouco visível. Um pipeline novo, uma integração com um sistema de origem ou uma
melhoria de desempenho na plataforma raramente aparecem numa reunião de diretoria, e por isso
são difíceis de priorizar, de explicar e de valorizar. Este capítulo mostra como o modelo de
Scrum com Kanban se aplica a esse tipo de trabalho: como escrever histórias quando o usuário
é outro sistema, onde colocar o débito técnico e como conciliar projetos novos com a
sustentação do que já existe.

\section{Histórias quando o usuário é um sistema}
\label{sec:historias-tecnicas}

O formato tradicional das histórias de usuário (\enquote{como tal pessoa, quero tal coisa,
para tal objetivo}) funciona mal quando o trabalho é \enquote{ingerir os dados do novo
sistema de pedidos}. Não há uma pessoa esperando na ponta, e a tentação é escrever histórias
puramente técnicas, como \enquote{criar tabela de pedidos na camada bronze}, que ninguém de
fora do time entende nem consegue priorizar.

Duas práticas ajudam. A primeira é amarrar o trabalho técnico ao valor que ele destrava, mesmo
que esse valor esteja algumas histórias à frente. \enquote{Ingerir os pedidos do novo sistema
para que o painel de vendas continue atualizado depois da migração} diz por que o trabalho
importa e quando ele precisa estar pronto. Quando o usuário é de fato outro sistema, a história
pode ser escrita do ponto de vista dele, com o contrato de dados como critério de aceite
(\veja{sec:contratos}).

A segunda é cuidar do tamanho. Uma boa história é pequena o bastante para atravessar o quadro
em poucos dias e independente o bastante para ser entregue sozinha \parencite{cohn2004}.
Projetos de engenharia tendem a gerar histórias enormes (\enquote{migrar a plataforma}) que
ficam semanas na coluna \enquote{fazendo}. Quebrá-las em fatias que entregam algo verificável,
como uma fonte migrada por vez, com os testes comparando o resultado antigo com o novo, torna
o progresso visível e o risco menor.

Os critérios de aceite de histórias técnicas também podem, e devem, virar testes escritos antes
(\veja{sec:tdd}): o volume carregado bate com a origem, a atualização acontece dentro do
prazo, nenhuma chave fica duplicada. Para engenharia, essa é a forma mais concreta de definir
o que é \enquote{pronto}.

\section{Débito técnico}
\label{sec:debito-tecnico}

A metáfora do débito técnico foi proposta por Ward Cunningham, em 1992, para explicar a
pessoas de negócio por que um código precisava ser revisto: entregar algo rápido, com uma
solução provisória, é como pegar um empréstimo, que acelera hoje e cobra juros depois, até que
seja pago com a reescrita \parencite{cunningham1992}. Em dados, os juros aparecem como
pipelines frágeis que quebram toda semana, modelos que ninguém entende e consultas lentas que
encarecem cada execução.

O débito técnico em si não é o problema. Às vezes faz sentido assumi-lo, de forma consciente,
para cumprir um prazo de negócio. O problema é o débito invisível, que ninguém registrou e que
só aparece na forma de incidentes. A prática mais útil é simples: todo débito assumido vira um
item no quadro, com a descrição do que foi feito de forma provisória e do risco que isso traz.
Assim, ele pode ser priorizado como qualquer outro trabalho.

Para que o débito seja pago de fato, muitos times reservam uma parte fixa da capacidade de cada
iteração para reduzi-lo \pendente{confirmar se eu usava uma porcentagem fixa e qual}. Outra
forma é aproveitar cada história que toca um modelo problemático para melhorá-lo um pouco, em
vez de deixar a melhoria para um projeto de reescrita que nunca é aprovado. E a IA ajuda
aqui: refatorar, documentar e acrescentar testes a código antigo é uma das tarefas em que os
agentes de programação mais economizam tempo, desde que alguém revise.

\section{Projeto e sustentação}
\label{sec:projeto-sustentacao}

Todo time de engenharia de dados divide a capacidade entre construir coisas novas e manter o
que já existe. A sustentação (falhas de carga, números estranhos, pedidos de ajuste, dúvidas
de usuários) chega todos os dias, em quantidade imprevisível, e tem prioridade, porque alguém
está decidindo com aqueles dados agora. Se ela compete diretamente com os projetos, os projetos
atrasam sempre, e o time vive a sensação de nunca terminar nada.

Há duas formas principais de organizar essa convivência. A primeira são duas raias no mesmo
quadro, com uma parte da capacidade reservada para a sustentação
(Figura~\ref{fig:quadro-scrumban}). A segunda é um rodízio: a cada período, uma pessoa do time
fica responsável pela sustentação, como um plantão, e o resto do time fica protegido das
interrupções. Os dois modelos funcionam; a escolha depende do tamanho do time e do volume de
sustentação. Separar em dois times diferentes também é possível em estruturas grandes, mas cria
o risco de que quem constrói nunca sinta o custo de manter o que construiu.

\rascunho{
  Confirmar qual dessas formas eu usava (raias, rodízio, outra) e por quê.
}

Uma situação recorrente merece atenção à parte: os projetos paralelos puxados por um projeto de
negócio, como o de master data do capítulo~\ref{cap:master-data}
(\veja{sec:master-data-puxado}). Eles entram no quadro como um épico próprio, com histórias
amarradas às do projeto principal, e não como tarefas escondidas dentro dele. Isso dá
visibilidade ao trabalho, permite medir quanto ele custou e, principalmente, deixa claro para o
negócio que o domínio entregue vai ser reaproveitado pelos próximos projetos.

Por fim, vale medir a sustentação. Quanto da capacidade do time ela consome, que tipo de
problema mais aparece, quais produtos geram mais chamados. Esses números mostram onde o débito
técnico está cobrando juros mais altos e ajudam a justificar o investimento em pagá-lo.

\rascunho{
  Relato de um projeto real de engenharia de dados, sem nomes de empresa ou cliente: como foi
  quebrado em histórias, que débito foi assumido e como a sustentação foi organizada durante o
  projeto.
}

\begin{leituras}
  \leitura{cohn2004}{Para escrever histórias pequenas e independentes, inclusive quando o
    trabalho é técnico.}
  \leitura{reis2022}{Para a visão completa do ciclo de vida da engenharia de dados, útil para
    planejar projetos de plataforma.}
  \leitura{forsgren2018}{Para medir a sustentação e a recuperação de falhas com os mesmos
    indicadores usados pelos melhores times de software.}
\end{leituras}
